\section{Integration and reproducibility}
\label{app:reproduction}

\subsection{Archive contents and source identity}

The archive is a self-contained research and integration bundle. Its
\code{article/} directory contains this PDF, the complete modular \LaTeX{}
source, bibliography, generated tables, and vector and raster figures.
\code{integration.patch} applies the implementation and its new tests to
the pinned repository. The matching \code{repo_overlay/} contains only
new or changed files, so it can be reviewed without extracting an entire
second repository over a working checkout.

The runnable final snapshot is at
\begin{lstlisting}
snapshots/final/Topology/UnknotRecognition/fast/
\end{lstlisting}
It contains the complete Python package, maintained tests, and required
fixtures. The neighboring \code{reports/} subdirectories contain historical
reference implementations imported by those tests. The archive omits the
unrelated large \code{fast/results/} collection and generated Python bytecode.
It includes complete predecessor and eager-prototype Python package snapshots
at \code{snapshots/baseline/} and \code{snapshots/eager/} for isolated comparison.
The final snapshot has the exact source bytes used by the guarded records.

The \code{experiments/} directory contains the independent worker harness,
summary and figure generators, completed test ledger, exhaustive local
review, and four principal records:
\begin{center}
\small
\begin{tabular}{lp{0.57\linewidth}}
\toprule
Record & Evidence retained \\
\midrule
\code{audit-guarded.json} & Source cases, all legacy comparisons, independent
replay outcomes, positive certificates, and bounded nondecisions. \\
\code{stages-guarded.json} & Full completed-call timings and structural
counters on the five stabilized-circle sizes. \\
\code{benchmark-guarded.json} & Every round and A/A arm for all 19 ordinary
whole-recognizer cases. \\
\code{audit-eager.json} & The rejected partial-batch policy, including the
Gordian work-limit regression. \\
\bottomrule
\end{tabular}
\end{center}

The original ordinary-corpus record is retained byte-for-byte as
\code{experiments/corpus.json}. Its SHA-256 is
\begin{lstlisting}
ed7ae3116c89a294e0bcf716cf027f4c80b6bcf48523e6a0bb0e4ecb97d7e89b
\end{lstlisting}
It comes from report 46's cyclic-overlap research overlay at the pinned
commit. The harness reads its source diagrams and expected verdicts; it
does not import that report's algorithms to obtain the new results.

The manifest records the pinned commit, file hashes, snapshot roles, and
verification results. Checksums identify source and evidence bytes; they
are provenance aids, not a replacement for mathematical or source replay
verification. The included verification script checks the manifest,
snapshot-to-experiment identities, and the exact corpus and harness hashes.

\subsection{Applying the implementation}

From a clean ProveIt checkout at the pinned commit, inspect and apply the
patch using the archive's absolute path:
\begin{lstlisting}[language=bash]
git apply --check /path/to/unknot_singleton_dag_20261009/integration.patch
git apply /path/to/unknot_singleton_dag_20261009/integration.patch
\end{lstlisting}
The delivered patch changes five existing Python modules, adds the
producer and independent verifier, adds two test modules, and documents
the option in \code{SINGLETON_DAG.md}. The patch does not modify the Rust
implementation. An integration onto a later commit requires a normal
review of intervening changes; a successful application against the pinned
base does not guarantee compatibility with arbitrary future revisions.

For the research writeup, the \code{article/} and \code{experiments/}
directories can be placed together in a dated research directory under
\code{Topology/UnknotRecognition}. Their relative layout is used by the
summary and figure generators. Preserve the root manifest and the complete
archive when exact offline reproduction of the predecessor comparison is
desired. Nothing in this bundle has been committed or pushed to the remote
repository.

\subsection{Replaying the tests and experiments}

The core package and research harness require Python 3.10 or later and use
the standard library. The recorded run used Python 3.12.14. The optional
normal-surface tests require Regina; version 7.4.1 was installed for the
complete no-skip test run. Figure regeneration uses Matplotlib. The
\LaTeX{} build requires a standard distribution containing the packages
listed in the document preamble, together with BibTeX.

From the archive root, use the supplied scripts:
\begin{lstlisting}[language=bash]
python3 scripts/verify_bundle.py
bash scripts/reproduce.sh tests
bash scripts/reproduce.sh audit
bash scripts/reproduce.sh stages
bash scripts/reproduce.sh benchmark
bash scripts/reproduce.sh eager
make -C article
\end{lstlisting}
The test command deliberately runs with the final \code{fast/} directory
as its working directory. Some maintained tests launch
\code{python -m fastunknot} as a subprocess; merely changing the parent
process's import path would not select the intended package for those calls.
The suite's completion ledger must report 1,008 tests run, not merely a
successful shell exit or a partially written test log.

The experiment commands write fresh results beneath \code{reproduced/}
and leave the original evidence untouched. They pass explicit predecessor,
current, corpus, snapshot, and output paths to the frozen harness. The
stage command uses the five published sizes and five measured rounds.
The eager command selects the retained eager source and uses a separate
snapshot directory. Run timing commands sequentially, without concurrent
CPU-intensive work. Elapsed times will vary with the environment; the
source identities, exact certificates, coverage, and structural mechanisms
are the reproducible non-timing claims.

The independent exhaustive review can also be run directly:
\begin{lstlisting}[language=bash]
python3 experiments/review_singleton_verifier.py \
  snapshots/final/Topology/UnknotRecognition/fast
\end{lstlisting}
Its expected output reports 3,698 cases: 360 accepted and 3,338 rejected.
The result includes the verifier's source hash and a check of the skipped
partial-dispatch handoff.

The original tables are regenerated from the supplied result records by
\begin{lstlisting}[language=bash]
python3 experiments/summarize_singleton_dag.py
python3 experiments/plot_results.py
\end{lstlisting}
The former script intentionally reads the neighboring named evidence files.
For a new experimental run, copy the summary script beside the fresh records
and give it a separate \code{--article-generated} destination, rather than
overwriting the published evidence or presenting new timings as the old run.

\subsection{A seed-metadata qualification}

The frozen audit record's top-level \code{seed} field contains the harness
master constant 261009817. Within \code{audit()}, the actual random-braid
generator uses the explicit constant 261008504. The timing generators use
261009818 for the ordinary workload and 261009819 for the stage experiment.
This distinction is visible in the preserved hashed harness; the complete
PD arrays are also stored in the audit. We retain the original record and
state the qualification here, instead of silently rewriting its metadata.
The 81 source cases include repeated small diagrams and represent 75
distinct raw PD arrays.

\subsection{What would invalidate a performance comparison}

Changing the source snapshot during a run, mixing package versions through
absolute imports, excluding certificate replay from only one arm, dropping
bounded nondecisions, or selecting only favorable input cases would change
the meaning of the comparison. The harness and stored records make those
choices inspectable. Future changes should publish their new hashes,
completed verdicts, all sample orders, resource allowances, and the same
stage-versus-whole-recognizer distinction before interpreting timing ratios.
